\section{Discussion}
\label{sec:discussion}

\paragraph{Design decisions.} Besides the selection through Wikidata and the
two-layer ontology (Sections~\ref{sec:collection} and~\ref{sec:ontology}),
inferred triples are \emph{materialised and stored separately}, so clients get
DBpedia types without a reasoner and the interface can mark what is inferred;
and \emph{one Python process with an in-memory rdflib graph} serves the
endpoint, the Linked Data routes and the interface. The 131 thousand triples
load in a few seconds; a triple store such as Virtuoso pays off when the data
grows.

\paragraph{Limitations.} The current system has six limitations.
\begin{itemize}
  \item \emph{Coverage.} 990 entities of one domain, not all of Vietnamese
        Wikipedia; each new class needs a seed query, ontology terms, an
        infobox mapping and a build function.
  \item \emph{Resolvable IRIs.} We do not control \term{vi.dbpedia.org}, so
        the IRIs dereference only through a running project server; a domain
        of our own or \term{w3id.org} would make them resolvable on the Web.
  \item \emph{Link drift.} 44 of the 777 English links do not reach a DBpedia
        article (Section~\ref{sec:links}); following
        \term{dbo:wikiPageRedirects} would fix twelve of them.
  \item \emph{Data gaps.} 129 career stations have no team IRI, 26
        universities have no province, and page links
        (\term{dbo:wikiPageWikiLink}) and \term{skos:broader} are not extracted.
  \item \emph{Serving.} The SPARQL endpoint has no query time limit, because
        rdflib cannot cancel a running query; a public deployment needs a
        stoppable process or a triple store that enforces a limit.
  \item \emph{Question answering.} Its quality depends on the language
        model and was not measured on a benchmark; the evidence graph is
        empty when the rows hold one end of each relation.
\end{itemize}
